\documentclass[11pt,a4paper]{article}

\usepackage[margin=1in]{geometry}
\usepackage{booktabs}
\usepackage{multirow}
\usepackage{hyperref}
\usepackage{xcolor}
\usepackage{amsmath}
\usepackage{natbib}
\usepackage{array}

\hypersetup{
  colorlinks=true,
  linkcolor=blue!60!black,
  citecolor=blue!60!black,
  urlcolor=blue!60!black
}

\title{%
  Correct Verdicts, Wrong Field: Attribution and Capability \\
  Are Dissociable When Small Language Models Follow a Rubric%
}

\author{Sophie D.\ Neilson \\
  \small Independent Researcher \\
  \small ORCID: \href{https://orcid.org/0009-0001-2430-1743}{0009-0001-2430-1743}
}

\date{September 2026}

\begin{document}
\maketitle

\begin{abstract}
A behavioral format assay changes the structure of a prompt and reads whether a
model's answers improve. Verdict accuracy is the usual measure. We show that it
is not sufficient. We gave three local language models a structured decision
rubric that names an explicit decision field, and asked each for a verdict and
for the field that determined it. Across six decision classes and two scenario
polarities, verdict accuracy sat at or near ceiling (0.96--1.00), while the share
of correct verdicts that cited the intended decision field ranged from 0.05 to
0.66. On the firing scenarios, only 27\% of correct verdicts cited the intended
field; the rest routed through a recognition axis the rubric does not designate
as the decision route. One 32B model answered correctly in every run while
naming the intended field in five of ninety-six. A grader scoring only verdicts
would credit the rubric's decision field for a result the model attributes
elsewhere. We then removed the two recognition axes the models were routing
through. Accuracy did not fall. The models re-routed to the intended field and
stayed correct: citation of the intended field rose from 63 to 96 runs for the
thinking model and from 5 to 65 for the 32B model. The shortcut is a preference,
not a crutch. The intended path was available throughout. Attribution is
dissociable from capability: the field a model reports is not the only field it
can use, and neither the verdict nor the self-report is a reliable proxy for the
other. We argue that a format assay must track the reasoning route, not only the
verdict, and that a model's self-reported route must itself be checked against an
intervention rather than trusted.
\end{abstract}

\section{Introduction}

A common way to test whether a change to a prompt's structure helps a language
model is to score the model's answers before and after. If accuracy rises, the
structure gets the credit. This paper is about a failure of that inference. A
model can reach the correct answer through a part of the prompt that the change
was not about, so that a rise in accuracy is not evidence that the intended
structure did any work.

We study this with a structured decision rubric. The rubric describes a class of
decision an agent faces, and it names one field, a scope condition, as the
field on which the verdict turns. It also carries two recognition axes that
describe what a violation and a correct action look like. The rubric states the
order of evaluation: read the scope condition, then render the verdict. We ask a
model for a verdict on a concrete scenario and, on the same turn, for the name of
the field that determined its verdict. The verdict is scored against ground
truth. The named field is scored against the field the rubric designates.

Two questions follow. First, does verdict accuracy track engagement with the
intended field? We find it does not: accuracy is near ceiling while intended-field
citation is far lower, and the gap is largest on the scenarios where the decision
class fires. Second, is a correct-but-off-route answer fragile: does it depend
on the field the model cites? We test this by removing the recognition axes the
models route through. It is not fragile. Accuracy holds and the models re-route
to the intended field. The route a model names is a preference, not a
requirement.

We call the resulting distinction \emph{attribution versus capability}. What a
model reports as the basis of its answer (attribution) is separable both from
whether the answer is correct (verdict accuracy) and from which fields the model
can actually use to be correct (capability). A behavioral format assay that reads
only the verdict sees none of this. It cannot tell a genuine effect of the
intended structure from a rise carried by an unintended field, and it cannot tell
a load-bearing field from one the model merely names.

\subsection{Related Work}

That a model's stated reasoning need not reflect the cause of its answer is
established for chain-of-thought prompting. \citet{turpin2023} show that models
give plausible reasoning that does not mention the features actually driving their
predictions, and \citet{lanham2023} measure how far a stated chain of thought is
load-bearing for the final answer. \citet{chen2025reasoning} report the same for
reasoning models, whose stated chains omit decisive influences such as an injected
hint. Our field-source measure is a compact form of
the same question for a rubric-following task: the model names the field it used,
and we check that name against an intervention rather than taking it at face
value. The distinction between an answer's correctness and the process that
produced it also motivates process-level supervision, where a reward is placed on
the reasoning steps rather than the outcome \citep{lightman2023}. Our result is a
reason that outcome scoring alone is insufficient for a format study, framed as a
measurement problem rather than a training one.

A model reaching the right answer through an unintended feature is the
shortcut-learning pattern \citep{geirhos2020}, described there for learned
representations; we observe it at the level of a model's self-reported field on a
structured prompt, and we add that the shortcut here is not load-bearing: the
intended route survives its removal. The rubric we use comes from a research
program on descriptive decision-point prompts \citep{neilson2026capc}; two adjacent
studies of format-dependent behavioral shifts in language models
\citep{wang2026narrative, liu2026behavioral} vary prompt structure and read
behavioral output, but neither separates the model's self-reported route from its
verdict, which is the measure this paper turns on.

\section{Materials and Method}

\subsection{The Rubric and the Field-Source Measure}

Each decision class is described by a rubric with a fixed shape. A \textbf{Y --
Decision terrain} block names the evaluation order and three parts: a
\textbf{Pull character} (why the locally-wrong action is attractive, present in
both firing and non-firing cases), a \textbf{Scope-operative-when} field (the
condition under which the decision class fires, with a verification procedure),
and a \textbf{Scope-not-operative-when} field (the condition under which it
does not, stated as a positive discriminator). Below the terrain sit two axes: a
\textbf{Marker axis} (an invariant describing what the violation looks like) and
an \textbf{Aim axis} (an invariant describing correct navigation). A \textbf{Rest
axis} describes territory where the class does not arise. The terrain block states:
``Evaluate in sequence: pull character $\rightarrow$ scope $\rightarrow$ verdict.
Do not render a verdict before completing the scope check.'' The Scope field is
therefore the field the rubric designates as the decision route; the Marker and
Aim axes are recognition and characterization, not the route.

We use six decision classes (Section~\ref{sec:classes}). Each has two scenarios: a
\emph{firing} scenario whose ground truth is ``yes'' (the class applies) and a
\emph{carve-out} scenario whose ground truth is ``no'' (a named exception holds).
The intended field for a firing scenario is Scope-operative-when; for a carve-out
scenario it is Scope-not-operative-when. Both are the Scope field.

The instrument is the response format. The model is asked to answer in four lines:
a \texttt{VERDICT} (yes/no), a \texttt{FIELD SOURCE} (the name of the field that
determined the verdict), an \texttt{OBSERVABLE} (the condition it cites), and the
\texttt{EVIDENCE} (the trace element it rests on). The verdict and the field
source are parsed deterministically. A \emph{wrong-path success} is a correct
verdict whose cited field is not the Scope field.

\subsection{Model and Sampling}\label{sec:sampling}

Three subjects, all Q4\_K\_M GGUF quantizations served by llama.cpp
(\texttt{llama-server}, build \texttt{b9445-af6528e6d}) over the raw
completion endpoint, so no server-side chat template sits between the
declared prompt and the model. Each model's instruct wrapper is published verbatim
(Appendix~\ref{app:materials}).

\begin{itemize}
  \item \textbf{Mistral-7B-Instruct-v0.3} (7.2B), non-thinking.
  \item \textbf{Qwen2.5-32B-Instruct} (32B), non-thinking.
  \item \textbf{Qwen3.8-27B} (27B), thinking enabled. Its reasoning trace is
    recorded per run.
\end{itemize}

Claude-family models are excluded as subjects: the rubric corpus is in their
training data. The sampler chain is declared complete and identical across
subjects: temperature 0.8, \texttt{min\_p} 0.05 as the sole truncation stage,
every other truncation stage and every penalty disabled, per-run seeds 1--8.
Penalties are off by design: the rubric induces repeated structure (axis-by-axis
traversal), and a repetition penalty would attenuate the behavior under study. The
served context was 32768 tokens for Mistral and Qwen3.8. Qwen2.5-32B was served at
8192 tokens: its weights plus the KV cache at 32768 exceed the 24\,GB GPU. The
prompts are about 3800 tokens with at most 1024 generated, so 8192 has ample
headroom and no run truncated; the served context is recorded per run and bounds
only the unused tail.

\subsection{Conditions and Arms}

Every run assembles one prompt: the shared rubric-definition document, then the
decision-class terrain, then the scenario trace, then the four-line instruction.
The study has two arms that differ only in the terrain:

\begin{itemize}
  \item \textbf{Base arm.} The full terrain, with the Marker and Aim axes
    present.
  \item \textbf{Ablation arm.} The same terrain with the Marker and Aim axes
    removed. The Y-Decision block (including both Scope fields), the Pull
    character, and the Rest axis are byte-identical to the base terrain; a
    per-file diff confirms only the two axes changed.
\end{itemize}

The design is 3 subjects $\times$ 2 arms $\times$ 6 classes $\times$ 2 scenarios
$\times$ 8 runs $=$ 576 runs. No run truncated, and the served model matched the
declared model in every run.

\subsection{Scoring}\label{sec:scoring}

The headline measures are deterministic. The verdict is parsed and compared to
ground truth. The field source is parsed and classified by the first field-family
keyword it names: Scope, Marker, Aim, Pull, or other. For each cell (subject,
arm, class, scenario) we report verdict accuracy over the eight runs, and among
the correct verdicts the share that cited the Scope field. A run's answer is the
four-line response; for the thinking model the reasoning trace precedes it and is
recorded separately, so the four parsed lines are the same shape for all subjects.
Read cells and direction, not exact counts: at eight runs a single-integer
difference is inside the noise band.

\subsection{Agents and Models}\label{sec:agents}

Table~\ref{tab:agents} lists every model and agent with a role in producing the
results. The three subjects are the models under test. A Claude session
authored the study and the analysis code; it is an agent, not a person, and it
judged nothing in the deterministic layer, which is a parser.

\begin{table}[h]
\centering
\small
\caption{Agents and models.}
\label{tab:agents}
\begin{tabular}{@{}>{\raggedright\arraybackslash}p{3.1cm}>{\raggedright\arraybackslash}p{3.6cm}>{\raggedright\arraybackslash}p{4.6cm}>{\raggedright\arraybackslash}p{2.4cm}@{}}
\toprule
Role & Model & Version and runtime & Date \\
\midrule
Subject (non-thinking) & Mistral-7B-Instruct-v0.3 & Q4\_K\_M, llama.cpp b9445, ctx 32768 & 2026-09-11 \\
Subject (non-thinking) & Qwen2.5-32B-Instruct & Q4\_K\_M, llama.cpp b9445, ctx 8192 & 2026-09-11 \\
Subject (thinking) & Qwen3.8-27B & Q4\_K\_M, llama.cpp b9445, ctx 32768, thinking on & 2026-09-11 \\
Study and analysis author & Claude (agent) & Version not recorded & 2026-09-11 \\
\bottomrule
\end{tabular}
\end{table}

\section{Results}

\subsection{Verdict Accuracy Overstates Field Engagement}

Verdict accuracy is at or near ceiling for every subject, while the share of
correct verdicts that cite the intended Scope field is far lower
(Table~\ref{tab:h1}). The models are right, and they largely say a different field
made them right.

\begin{table}[h]
\centering
\small
\caption{Base arm, per subject, over 96 runs each (12 cells $\times$ 8).}
\label{tab:h1}
\begin{tabular}{@{}lccc@{}}
\toprule
Subject & Verdict accuracy & Scope-cited / correct & Wrong-path / correct \\
\midrule
Mistral-7B & 0.96 (92/96) & 0.66 & 0.34 \\
Qwen2.5-32B & 1.00 (96/96) & 0.05 & 0.95 \\
Qwen3.8 & 1.00 (96/96) & 0.66 & 0.34 \\
\bottomrule
\end{tabular}
\end{table}

Qwen2.5-32B is the sharpest case. It answered correctly in all 96 runs and named
the Marker axis in 88 of them, the Scope field in five. A grader scoring only its
verdicts would read a perfect score and credit the scope structure for a result
the model attributes almost entirely to a recognition axis.

The gap is widest on the firing scenarios. Pooling the base arm across subjects,
verdict accuracy is 0.99 on both polarities (142/144 each), but intended-field
citation among correct verdicts is 0.27 on the firing scenarios against 0.64 on
the carve-outs. When the class applies, roughly three of four correct ``yes''
verdicts are reached by naming the Marker axis's firing condition rather than the
scope condition the rubric puts on the decision path.

\subsection{Removing the Shortcut Does Not Break the Answer}

If a correct verdict depended on the recognition axis the model cited, removing
that axis should collapse the cell. It does not. In the ablation arm, with the
Marker and Aim axes deleted, verdict accuracy holds at ceiling for all three subjects.
The base-versus-ablation accuracy differential between shortcut-routing cells and
scope-routing cells (classified from the base arm) is within noise of zero for
every subject; the only movements are two small \emph{improvements} under ablation
for Mistral (one carve-out cell 6/8 $\rightarrow$ 8/8, one firing cell 6/8
$\rightarrow$ 7/8).

What changes is the cited field. With the recognition axes gone, the models route
to the Scope field and stay correct (Table~\ref{tab:h2}). The thinking model cites
the Scope field in all 96 ablation runs, up from 63; the 32B model rises from 5 to
65.

\begin{table}[h]
\centering
\small
\caption{Field-source class over all 96 runs per subject per arm. Verdict
accuracy is unchanged between arms (near ceiling throughout).}
\label{tab:h2}
\begin{tabular}{@{}lll@{}}
\toprule
Subject & Base arm & Ablation arm \\
\midrule
Mistral-7B & scope 63, pull 15, marker 9, aim 1, other 8 & scope 73, pull 21, other 2 \\
Qwen2.5-32B & marker 88, scope 5, other 3 & scope 65, other 31 \\
Qwen3.8 & marker 32, scope 63, other 1 & scope 96 \\
\bottomrule
\end{tabular}
\end{table}

The intended path was available throughout. The models preferred the recognition
axis when it was present and used the scope condition when it was not, at no cost
to accuracy. The shortcut is a preference, not a crutch.

\section{Discussion}

\subsection{Attribution Is Not Capability}

Three quantities that a verdict-only assay silently conflates come apart here. The
\emph{verdict} is whether the answer is correct. The \emph{attribution} is the
field the model names as its basis. The \emph{capability} is which fields the
model can in fact use to be correct. Verdict accuracy is near ceiling regardless
of the other two. Attribution is mostly off the intended field in the base arm.
Capability includes the intended field the whole time, as the ablation shows.
Knowing the verdict tells you nothing about the attribution, and knowing the
attribution tells you nothing about the capability.

For a behavioral format assay this is a direct constraint. A rise in verdict
accuracy after a structural change is not evidence that the intended structure did
the work: the model may reach the answer through another field, as ours did three
times in four on the firing scenarios. And a field the model never names may still
be one it can use, so the absence of a citation is not evidence of an absent
capability. The assay must track the route, and it must not read the model's
self-reported route as ground truth.

\subsection{The Self-Report Must Be Checked, Not Trusted}

The field-source line is the model's account of its own decision. Taken at face
value it would say that Qwen2.5-32B is a Marker-axis reasoner that essentially
ignores the scope condition. The ablation refutes that: remove the Marker axis and
the same model cites the scope condition and stays correct. The self-report
described a preference, not a dependency. This is the chain-of-thought faithfulness
problem \citep{turpin2023, lanham2023, chen2025reasoning} in miniature and with a cheap control: a
one-line self-report, plus one ablation that removes the named field, separates
what the model says it used from what it needs. We recommend the pairing (ask
for the route, then delete the route and re-measure) as a routine check for any
assay that reads a model's stated basis.

\subsection{Relation to the Orientation Study}\label{sec:orientation}

An earlier small study (four runs per cell, on a since-retired inference runtime
with an unrecorded sampler) had found that adding structure to one decision class
drove its firing scenario from three-of-four correct to zero, and read this as
evidence of no reliable correct path beneath a shortcut. The present study, run
clean and at larger scale, does not support that reading. The earlier
intervention \emph{added} a routing block; ours \emph{removes} the shortcut. When
the shortcut is removed, the answer survives. We take the earlier collapse to have
been an effect of the added structure, not evidence of a missing path, and we
report the correction here because the orientation study's framing motivated this
one.

\subsection{Limitations}\label{sec:limitations}

Eight runs per cell support statements about direction and about large gaps, not
about small differences between counts. The field-source classifier keys on the
first field-family term a response names; a minority of ablation-arm responses
(31 of Qwen2.5-32B's 96) name a field the classifier records as ``other'' rather
than resolving to a specific family, and we have not yet read each to its named
field. Those are counted as non-Scope, which does not affect the direction of any
result. The observable-and-evidence quality of each response (whether the cited
condition is the discriminating one) is a separate, judge-scored layer we do
not report here. The three subjects are small open-weight models on one GPU; we do
not claim the exact rates transfer to larger or hosted models, only that the
dissociation appears across two families, two scales, and both reasoning modes.
Qwen2.5-32B ran at a smaller served context than the other two subjects; the
prompts fit within it, so this bounds only unused context, but it is an
inconsistency in the rig, recorded rather than hidden.

\section{Conclusion}

A structured decision rubric that names its decision field lets us separate three
things a verdict-only assay treats as one: whether a model is right, which field
it says made it right, and which fields it can use to be right. Across three small
language models the first was near ceiling, the second was mostly off the intended
field, and the third included the intended field throughout. Verdict accuracy did
not track intended-field engagement, and the model's self-report described a
preference that an ablation showed was not a dependency. A behavioral format assay
should therefore measure the reasoning route, not only the verdict, and should
check a model's stated route against an intervention rather than trust it.
Attribution and capability are dissociable, and neither is read off the verdict.

\subsection{Data Availability}

The complete study (rubric-definition document, the six decision-class terrains
in both the full and ablated forms, the twelve scenarios, the four-line
instruction, the generated study definitions, every run's rendered prompt,
response, and (for the thinking model) reasoning trace, the scored grids, and the
scoring script) is available at \url{https://github.com/sophdn/corpos-lab}
under the \texttt{studies/wrong-path-field-source/} directory.

\bibliographystyle{plainnat}

\begin{thebibliography}{9}

\bibitem[Chen et~al.(2025)]{chen2025reasoning}
Chen, Y., Benton, J., Radhakrishnan, A., Uesato, J., Denison, C., Schulman, J.,
  Somani, A., Hase, P., Wagner, M., Roger, F., Mikulik, V., Bowman, S.~R.,
  Leike, J., Kaplan, J., and Perez, E. (2025).
\newblock Reasoning Models Don't Always Say What They Think.
\newblock \emph{arXiv preprint arXiv:2505.05410}.

\bibitem[Geirhos et~al.(2020)]{geirhos2020}
Geirhos, R., Jacobsen, J.-H., Michaelis, C., Zemel, R., Brendel, W., Bethge, M.,
  and Wichmann, F.~A. (2020).
\newblock Shortcut learning in deep neural networks.
\newblock \emph{Nature Machine Intelligence}.
  \href{https://doi.org/10.1038/s42256-020-00257-z}{doi:10.1038/s42256-020-00257-z}.

\bibitem[Lanham et~al.(2023)]{lanham2023}
Lanham, T., Chen, A., Radhakrishnan, A., Steiner, B., Denison, C., Hernandez, D.,
  Li, D., Durmus, E., Hubinger, E., Kernion, J., Lukošiūtė, K., Nguyen, K.,
  Cheng, N., Joseph, N., Schiefer, N., Rausch, O., Larson, R., McCandlish, S.,
  Kundu, S., Kadavath, S., Yang, S., Henighan, T., Maxwell, T., Telleen-Lawton, T.,
  Hume, T., Hatfield-Dodds, Z., Kaplan, J., Brauner, J., Bowman, S.~R., and
  Perez, E. (2023).
\newblock Measuring Faithfulness in Chain-of-Thought Reasoning.
\newblock \emph{arXiv preprint arXiv:2307.13702}.

\bibitem[Lightman et~al.(2023)]{lightman2023}
Lightman, H., Kosaraju, V., Burda, Y., Edwards, H., Baker, B., Lee, T., Leike, J.,
  Schulman, J., Sutskever, I., and Cobbe, K. (2023).
\newblock Let's Verify Step by Step.
\newblock \emph{arXiv preprint arXiv:2305.20050}.

\bibitem[Liu et~al.(2026)]{liu2026behavioral}
Liu, H., Liu, R., Xu, H., Han, J., Fang, S., Yan, S., Deng, H., Wang, G., and
  Zou, N. (2026).
\newblock Your LLM, Your Style: Behavioral Mode Axes for LLM Behavioral Control.
\newblock \emph{arXiv preprint arXiv:2608.10703}.

\bibitem[Neilson(in preparation)]{neilson2026capc}
Neilson, S.~D. (in preparation).
\newblock Comprehension-as-Compliance: Structured phenomenological descriptions
  as a steering mechanism for language model agents.

\bibitem[Turpin et~al.(2023)]{turpin2023}
Turpin, M., Michael, J., Perez, E., and Bowman, S.~R. (2023).
\newblock Language Models Don't Always Say What They Think: Unfaithful
  Explanations in Chain-of-Thought Prompting.
\newblock \emph{arXiv preprint arXiv:2305.04388}.

\bibitem[Wang et~al.(2026)]{wang2026narrative}
Wang, Y., Lester, J., and Srivastava, S. (2026).
\newblock The Story Shapes the Agent: Narrative Priors in LLM Behavior.
\newblock \emph{arXiv preprint arXiv:2607.18566}.

\end{thebibliography}

\appendix

\section{Materials}\label{app:materials}

The instruction, given to every subject after the assembled material, is
reproduced verbatim:

\begin{quote}
Read the glyph and the trace above.

Respond using exactly this format --- one line per field, no additional text:

VERDICT: yes or no \\
FIELD SOURCE: the name of the specific field in the glyph specification that
determined your verdict \\
OBSERVABLE: quote or closely paraphrase the specific condition from that field \\
EVIDENCE: state which element of the trace your verdict rests on
\end{quote}

The per-subject instruct wrappers, into which the assembled material is
substituted, are: \texttt{[INST] \{prompt\} [/INST]} for Mistral;
\texttt{<|im\_start|>user$\backslash$n\{prompt\}<|im\_end|>$\backslash$n<|im\_start|>assistant$\backslash$n}
for both Qwen subjects, with no pre-filled think block so Qwen3.8 reasons.

\section{Decision Classes}\label{sec:classes}

The six decision classes are \texttt{cas} (companion-artifact-scope-gap),
\texttt{cgu} (conditional-gate-uniform-default), \texttt{fsb}
(formal-step-context-bypass), \texttt{gop} (governed-operation-protocol-bypass),
\texttt{psc} (parent-state-check-bypass), and \texttt{scb}
(structural-ceiling-bypass). Each terrain, in full and ablated form, and each
scenario, is in the study directory named in Data Availability.

\end{document}
